\section{The finite stopped correlated ledger}
\label{sec:terminal-modal-finite-stopped-ledger}

The stopped leading comparison from
\cref{lem:terminal-modal-stopped-leading-baseline} survives at the named
positive $q$.  The proof must keep the fold: the derivative packet has a
larger stopped Abel gain on the fixed $2m$-cycle than on a growing path, and
the literal antipode is a single coordinate.

Put $\delta=1-q$ and $s=\delta^{m-1}$.  Let
$C_{m+1}^{\mathrm{act}}$ be the correction at literal full-face entry and
$C_{m+1}^{\mathrm{prop}}$ its hypothetical degree-two continuation.  The
entry data from \cref{lem:terminal-modal-directed-packet} and
\cref{lem:terminal-modal-static-remainder} obey the exact identities
\begin{equation}
 c=-C_{m+1}^{\mathrm{act}},\qquad
 d=\delta C_m-C_{m+1}^{\mathrm{act}}+r_m e_m,\qquad
 r_m=q\eta P_m-\frac{q^3}{5}.
 \label{eq:terminal-modal-finite-stopped-cd}
\end{equation}
Indeed $g=-G_{m+1}^{\rm id}$, so $c=r_0^x-g=-C_{m+1}^{\rm act}$.
Moreover exact velocity transport gives
$q r_0^v=A_{m+1}-\delta A_m+r_me_m$, while the directed
deconvolution packet is
$e=-G_{m+1}^{\rm id}+\delta G_m^{\rm id}$.  Since
$C_n=X_n-G_n^{\rm id}$ and $X_n=-A_n$, subtraction proves the second
identity.

Write
\begin{equation}
 \mathcal Q_k=-\mathcal H_kc-\mathcal T_kd,
 \qquad1\le k<m.
 \label{eq:terminal-modal-finite-stopped-Q}
\end{equation}
By \eqref{eq:terminal-modal-early-half-kernel} and
\eqref{eq:terminal-modal-half-velocity-static-identity}, this is exactly the
non-directed part of the scaled linear half-retention quantity.

\begin{lemma}[Folded proper-source perturbation; proved here]
\label{lem:terminal-modal-finite-stopped-proper-perturbation}
Relative to the leading stopped proper trajectory, the finite-$q$ initial,
derivative, and mass perturbations in $\mathcal Q_k/q^3$ have absolute value
strictly below, respectively,
\begin{equation}
 \frac3{5120},\qquad \frac{11}{720},\qquad \frac{43}{1200},
 \qquad1\le k<m.
 \label{eq:terminal-modal-finite-stopped-proper-losses}
\end{equation}
These bounds include the standard hypothetical degree-two source at $n=m$.
\end{lemma}

\begin{proof}
Use the rescaling
\eqref{eq:terminal-modal-finite-rescaling}.  The homogeneous recurrence is
the leading recurrence exactly, and the perturbation sources are
\begin{equation}
 h_n(1,2,-3)+\mu_n(0,0,1),\qquad
 h_n=\frac{\varepsilon_n}{4},\qquad2\le n\le m,
 \label{eq:terminal-modal-finite-stopped-sources}
\end{equation}
together with reflection.  The endpoint $n=m$ here is the standard source
in $C_{m+1}^{\mathrm{prop}}$, not the later literal degree-one defect.
The sixteen-cell proof of
\cref{lem:terminal-modal-final-source-variation} extends verbatim through
this source: its last parameter is $r=m-2$, and
$\chi^{m-1}>\chi^m\ge7/8$.  Hence
\begin{equation}
 \operatorname{TV}(\varepsilon_2,\ldots,\varepsilon_m)<\frac1{300},
 \qquad
 |\varepsilon_m|+\operatorname{TV}(\varepsilon_2,\ldots,\varepsilon_m)
 <\frac{11}{1800}.
 \label{eq:terminal-modal-finite-stopped-epsilon-TV}
\end{equation}

We first prove the folded Abel constant.  Fix a target $0\le j\le m$ and
put $N=m+k+1$.  For a source at $n$, its propagation time is
$h=N-n$, traversed in decreasing order as source chronology increases.
For the positive orientation, the line distance parameter is
\begin{equation}
 d_+=N-2-j\ge0.
\end{equation}
The lobe calculation in
\eqref{eq:terminal-modal-finite-derivative-negative-lobe}--
\eqref{eq:terminal-modal-finite-derivative-lobes} shows that every stopped
chronological prefix of this branch lies in $[-4,2]$.  Its possible
$-2m$ alias has parameter at most $-2$ and is zero.  The $+2m$ distance
alias lies entirely in a nonnegative lobe and contributes a stopped prefix
in $[0,2]$.  This term cannot be discarded: for example, at
$(m,k,N,j,n)=(8,7,16,4,2)$ it equals
$w_{14}(26)=15/8192>0$.

For the reflected orientation, the main distance is
$d_-=N-2+j$ and every propagation time obeys $h\le d_-$, so this branch is
inside its nonnegative lobe and lies in $[0,2]$.  Its only possible wrap has
parameter
\begin{equation}
 d_0=d_--2m.
\end{equation}
Indeed its $+2m$ alias vanishes: since $h\le N-2$ and $k<m$,
\begin{equation}
 d_-+2m-(2h-1)\ge 2m+j+3-N=m+j+2-k>0,
\end{equation}
and larger aliases are farther outside the support.  Thus $d_0$ is the only
reflected wrap that remains.
It is zero if $d_0\le-2$; if $d_0\ge0$, its stopped prefix lies in
$[-4,2]$.  In the remaining case $d_0=-1$, direct substitution in
\eqref{eq:terminal-modal-finite-l-coefficients} gives
\begin{equation}
 w_h(-1)=-\frac3{2^{h-1}},\qquad h\ge k+1,
\end{equation}
so the magnitude of any stopped tail is at most
$6/2^k\le3$.  Thus all two-orientation aliases together lie in $[-8,8]$;
the exceptional $d_0=-1$ case is smaller.

At $n=m$ the two orientations would count the antipodal coefficient $-3$
twice, whereas literal even extension counts it once.  The literal packet
is therefore the two-orientation packet plus $3e_m$.  This correction is
nonnegative after propagation and is at most $3/2$, by the folded
one-atom bound proved below.  It affects only the last chronological term,
so every literal prefix lies in $[-8,19/2]$.  A uniform folded Abel gain of
ten is consequently valid.  Applying Abel summation to
\eqref{eq:terminal-modal-finite-stopped-epsilon-TV} gives
\begin{equation}
 10\left(|h_m|+\operatorname{TV}(h_2,\ldots,h_m)\right)
 =\frac52\left(|\varepsilon_m|+\operatorname{TV}(\varepsilon_2,\ldots,
 \varepsilon_m)\right)<\frac{11}{720}.
 \label{eq:terminal-modal-finite-stopped-derivative-loss}
\end{equation}

For the mass and base packets we use a folded coefficient bound.  If
$2\le h\le2m-1$, then
\begin{equation}
 \ell_h(d)=\frac12\Pr\{\operatorname{Bin}(h-2,1/2)\ge d-1\},
 \qquad d\ge1,\qquad \ell_h(0)=\frac12.
\end{equation}
The two aliases of one cycle atom have distances $d$ and $2m-d$; their tail
thresholds sum to $2m-2>h-2$.  Binomial symmetry therefore makes the two
tail probabilities sum to at most one.  Hence every folded coefficient of
one cycle atom is at most $1/2$.  An interior even point mass has two cycle
atoms and coefficient at most one; the self-reflected source $n=m$ has only
one and is smaller.  Since
$|\mu_n|<43q/75$ for the $m-1$ sources $2\le n\le m$, their total loss is
strictly below
\begin{equation}
\frac{43q(m-1)}{75}<\frac{43}{1200}.
\end{equation}
The scalar identity
\eqref{eq:terminal-modal-finite-mu-M} is valid at $n=m$ for the hypothetical
degree-two continuation: it uses only $P_m$, $P_{m-1}$, and the fixed-prefix
frontier relation.  Thus the displayed mass bound extends to $\mu_m$
directly and is not imported from the literal degree-one endpoint formula.

Finally the exact initial perturbations have even-line norms $q/5$ and
$q-q^2/5<q$.  Here $N=m+k+1\le2m$, so the same alias bound applies to
$L_{N-1}^{\circ}$ and $L_{N-2}^{\circ}$; convolution by $B_0$ preserves the
maximum coefficient.  The initial loss is at most
$q/10+q/2=3q/5\le3/5120$.  Multiplication by the common
$s=\delta^{m-1}\le1$ can only reduce these three normalized losses.  This
proves \eqref{eq:terminal-modal-finite-stopped-proper-losses}.
\end{proof}

\begin{lemma}[Literal endpoint perturbation; proved here]
\label{lem:terminal-modal-finite-stopped-endpoint-perturbation}
Let $c_0=1/10-3p/8$ and $d_0=p/8-1/10$, where $p=P_m/q^2$.
The literal endpoint atoms obey
\begin{equation}
 \left|\frac{c_{\rm end}}{q^3}-sc_0\right|<\frac7{160},\qquad
 \left|\frac{d_{\rm end}}{q^3}-sd_0\right|<\frac1{40}.
 \label{eq:terminal-modal-finite-stopped-endpoint-errors}
\end{equation}
Their contribution to $\mathcal Q_k/q^3$ has absolute value below $3/128$
for every $1\le k<m$.
\end{lemma}

\begin{proof}
Put
\begin{equation}
 p_n=\frac{P_n}{q^2},\qquad S_n=\frac{\sigma_n}{q^3},\qquad
 \zeta_n=S_n-\left(\frac32p_n-\frac25\right).
\end{equation}
The exact frontier recurrence
\eqref{eq:terminal-modal-velocity-frontier-scaled} gives
\begin{align}
 \zeta_n&=\frac{\eta\zeta_{n-1}+Q_n}{1+q},
 \label{eq:terminal-modal-finite-stopped-zeta-recurrence}\\
 Q_n&=\eta(p_{n-1}-p_n)
 -q\left(\frac32+\frac q2\right)p_n+\frac{2q}{5}+\frac{q^2}{5}.
 \label{eq:terminal-modal-finite-stopped-Qn}
\end{align}
We record a direct Lipschitz bound for the frontier profile.  Differentiating
\eqref{eq:terminal-modal-frontier-optimum} with respect to
$t=\chi^n$ gives
\begin{equation}
 p'(t)=\frac4{1-q^2}
 \frac{1-t^2+2t/5}{(1+t^2)^2}<1
 \qquad(7/8\le t\le1).
\end{equation}
Indeed the three factors are bounded by $5$, $3/5$, and $1/3$,
respectively.  Also
$t_{n-1}-t_n<2q/(1-q)<3q$, so
\begin{equation}
 0<p_{n-1}-p_n<4q.
 \label{eq:terminal-modal-finite-stopped-p-difference}
\end{equation}
Using $p_n\le33/16$ and $q\le1/1024$ in
\eqref{eq:terminal-modal-finite-stopped-Qn} yields
\begin{equation}
 \frac{|Q_n|}{q}
 <2+\left(\frac32+\frac1{2048}\right)\frac{33}{16}
   +\frac25+\frac1{5120}
 =\frac{900293}{163840}<\frac{11}{2}.
 \label{eq:terminal-modal-finite-stopped-Qn-bound}
\end{equation}
Since $\eta/(1+q)=\delta/2$, while
$S_1=1/q-1/5$, one has
\begin{equation}
 \zeta_1=\frac1q+\frac15-\frac32p_1,\qquad0<\zeta_1<\frac1q.
\end{equation}
Indeed the lower inequality follows from
$p_1\le33/16$ and $1/q\ge1024$, while the upper follows from
$p_1\ge3488/1921>2/15$.  Iteration of
\eqref{eq:terminal-modal-finite-stopped-zeta-recurrence} now gives
\begin{equation}
 |\zeta_{m-1}|
 <\frac{1}{q\,2^{m-2}}+\frac{11q}{(1+q)^2}
 <\frac1{1024}+\frac{11}{1024}=\frac3{256}.
 \label{eq:terminal-modal-finite-stopped-zeta-bound}
\end{equation}
The first strict inequality in the last line uses
$16m/2^{m-2}<1/1024$ at $m=64$ and its decreasing ratio thereafter.

Let $p_-=p_{m-1}$.  Direct substitution in the standard and literal final
source formulas gives
\begin{align}
 W_{\rm std}&=\frac{q\eta^2P_{m-1}}{4(1+q)}
 -\frac{\delta qP_m}{4}+\frac{(3+q)q^3}{20},
 \label{eq:terminal-modal-finite-stopped-Wstd}\\
 W_{\rm end}&=-\frac{\delta\eta\sigma_{m-1}}4
 -\frac{q\delta P_m}{2}+\frac{q\eta\delta P_{m-1}}4+\frac{q^3}{5}.
 \label{eq:terminal-modal-finite-stopped-Wend}
\end{align}
Thus
\begin{align}
 c_{\rm end}&=W_{\rm end}-W_{\rm std}
 =\delta\left(-\frac{\eta\sigma_{m-1}}4-\frac{qP_m}{4}
 +\frac{q\eta P_{m-1}}8+\frac{q^3}{20}\right),
 \label{eq:terminal-modal-finite-stopped-cend}\\
 d_{\rm end}&=c_{\rm end}+q\eta P_m-\frac{q^3}{5}.
 \label{eq:terminal-modal-finite-stopped-dend}
\end{align}
Define
\begin{equation}
 E=\frac{p-p_-}{8}+q^2\left(\frac{p_-}{8}-\frac1{20}\right)
 -\frac{\eta\zeta_{m-1}}4.
\end{equation}
Using \eqref{eq:terminal-modal-finite-stopped-p-difference},
$p_-\le33/16$, and
\eqref{eq:terminal-modal-finite-stopped-zeta-bound} gives
\begin{equation}
 |E|<\frac q2+\frac{q^2}{4}+\frac3{2048}
 \le\frac{8193}{4194304}<\frac1{500}.
 \label{eq:terminal-modal-finite-stopped-E-bound}
\end{equation}

Let
$R_c=c_{\rm end}/q^3-sc_0$ and
$R_d=d_{\rm end}/q^3-sd_0$.  Exact simplification gives
\begin{align}
 R_c&=\delta E+(\delta-s)c_0,
 \label{eq:terminal-modal-finite-stopped-Rc}\\
 R_d&=\delta E+\frac p8\{\delta(1+4q)-s\}
 -\frac{q+1-s}{10}.
 \label{eq:terminal-modal-finite-stopped-Rd}
\end{align}
Now $1-s<1/16$, $p\le2$, and $|c_0|\le13/20$, so
\begin{equation}
 |R_c|<\frac1{500}+\frac{13}{320}
 =\frac{341}{8000}<\frac7{160}.
\end{equation}
For the second remainder put $x=1-s$.  Its non-$E$ part is
\begin{equation}
 x\left(\frac p8-\frac1{10}\right)
 +q\left(\frac{p(3-4q)}8-\frac1{10}\right),
\end{equation}
which is nonnegative and below
$3/320+13/20480=41/4096$.  Hence
\begin{equation}
 |R_d|<\frac1{500}+\frac{41}{4096}
 =\frac{6149}{512000}<\frac1{40}.
\end{equation}
This proves \eqref{eq:terminal-modal-finite-stopped-endpoint-errors}.

Finally, when $k<m$, the two directed supports in
\begin{equation}
 \mathcal H_k=\frac14\left(
 \mathcal S_+\mathcal A_+^{k-1}
 +\mathcal S_-\mathcal A_-^{k-1}\right)
\end{equation}
cannot both reach a fixed target from the antipode.  Its endpoint coefficient
is therefore at most $1/4$.  The endpoint coefficient of $\mathcal T_k$ is
at most $1/2$ by
\eqref{eq:terminal-modal-half-velocity-tail}, with no alias because $k<m$.
The total endpoint perturbation is consequently below
\begin{equation}
 \frac14\frac7{160}+\frac12\frac1{40}=\frac3{128}.
\end{equation}
\end{proof}

\begin{theorem}[Finite stopped correlated positivity; proved here]
\label{thm:terminal-modal-finite-stopped-positivity}
Let $m\ge64$.  For every $1\le k<m$ and every full-path row,
\begin{equation}
 \frac{\mathcal Q_k(j)}{q^3}>\frac{1807}{115200}>0.
 \label{eq:terminal-modal-finite-stopped-positivity}
\end{equation}
Consequently the early half-retention condition
\eqref{eq:terminal-modal-early-half-retention} holds unconditionally for
$qk<3/50$; indeed that inequality gives $k<24m/25<m$.  Projection and
envelope subtraction are inactive throughout
that early window.
\end{theorem}

\begin{proof}
The leading contribution is normalized by $q^3s$, not merely by $q^3$.
By \cref{lem:terminal-modal-stopped-leading-baseline}, it is at least
$31/320$ in that normalization.  Bernoulli's inequality gives
$s=\delta^{m-1}>15/16$, so in absolute $q^3$ units it is larger than
\begin{equation}
 \frac{15}{16}\frac{31}{320}=\frac{93}{1024}.
\end{equation}
Subtracting the three proper-source losses from
\cref{lem:terminal-modal-finite-stopped-proper-perturbation} and the endpoint
loss from
\cref{lem:terminal-modal-finite-stopped-endpoint-perturbation} leaves
\begin{equation}
 \frac{93}{1024}
 -\left(\frac{11}{720}+\frac3{5120}
       +\frac{43}{1200}+\frac3{128}\right)
 =\frac{1807}{115200}>0.
\end{equation}
The directed ideal difference in
\eqref{eq:terminal-modal-early-half-kernel} is nonnegative.  Therefore
\eqref{eq:terminal-modal-finite-stopped-positivity} proves strict linear
half-retention.  The implication in
\cref{lem:terminal-modal-early-half-interface} then proves the literal early
projection and envelope claims.
\end{proof}

This closes the early nonlinear interface without proving the stronger
separate inequality
\eqref{eq:terminal-modal-short-early-convolution-open}.  The static bound
\eqref{eq:terminal-modal-static-target-assumed}, needed by the late-envelope
lemma, is not proved by this stopped ledger alone.  It is closed for all
sufficiently large $m$ by the static synthesis in
\cref{thm:terminal-modal-static-asymptotic-closure}.
